\song{Bon soir mademoiselle Paris}
\artist{olympic}

\vers[1]{
\chord{Em}Mám v \chord{H\7}kapse jeden \chord{Em}frank, jsem \chord{D}nejbohatší z \chord{G}bank nad Sei\chord{H\7}nou.\\
\chord{Em}Mám \chord{H\7}víc než krupi\chord{Em}ér, \chord{D}stíny Sacre-\chord{G}Coeur\footnote{[sakré kér] = Nejsvětější Srdce, slavná pařížská bazilika na Montmartru} nade \chord{H\7}mnou.
}

\chorus{
\chord{C}Láska je \chord{D}úděl \chord{G}tvůj, \chord{Am}Pán Bůh tě \chord{H\7}opa\chord{Em}truj.\\
Bon so\chord{C}ir mademoi\chord{D}selle Pa\chord{Em}ris\footnote{[bonsuár madmoazel parí] = Dobrý večer, slečno Paříži.}, bon so\chord{C}ir mademoi\chord{D}selle Pa\chord{Em}ris.
}

\vers[2]{
Znám boulevard Saint-Michel\footnote{Významná ulice v Paříži na hranici 5. a 6. obvodu. Protože se nachází v blízkosti Sorbonny, stal se v květnu 1968 jedním z hlavních míst pouličních bojů mezi policií a studenty. Pařížské studentské nepokoje v květnu 1968 byly protesty proti konzervativní společnosti, autoritářství a zastaralému školskému systému, které se postupně rozšířily i mezi dělníky a vyústily v generální stávku. Hnutí požadovalo větší osobní svobodu a společenské změny a výrazně ovlivnilo francouzskou kulturu, politiku i vztah k autoritám.}, tam jsem včera šel s Marie Claire.\\
Vím, jak zní z úst krásných žen ta slůvka \uv{Car je t'aime, oh, ma chère.}\footnote{[kar ž’tem, ó, ma šér] = Protože tě miluji, ó, má drahá.}
} \cue{R}
